% Annex of auxiliary functions (specification only). Each entry states the function's meaning and the
% cases not shown; the cases are shared with the paper's figures. In a case, a side condition
% (if or where) goes in the condition column; a case too wide for that has its condition on a
% continuation line indented under the right-hand side.
\chapter{Auxiliary functions (normative)}
\label{ch:auxiliary-functions}

This annex is normative. It defines the auxiliary functions used by the well-formedness rules
(\secref{well-formedness}) and the operational semantics (\secref{operational-semantics}). When more than one
case applies, the first defines the value, so a function with overlapping cases is still determined: two
equivalent types join to the second of them, for instance.

\section{Syntactic analyses}
\label{sec:auxiliary-syntactic}

The syntactic analyses introduced in \secref{well-formedness} are defined by their principal cases; the remaining
cases are stated with each.

\begin{auxfunction}{\assignsS(t)}{set of variables}{aux:assigns}
Variables assigned anywhere in $t$, not descending into nested functions, which have their own scope. It is empty on the statements not shown and is the union over the branches of a
conditional, and on a module body $\seq{\iota}\;t$ acts on $t$.

\funccases{assigns}
\end{auxfunction}

\begin{auxfunction}{\binds(p), \binds(s)}{set of variables}{aux:binds}
Variables introduced by pattern $p$, or by the generators of qualifier sequence $\seq{g}$\Added{, or by the patterns
within statement $s$, not descending into nested functions}. On the cases
not shown it recurses structurally, taking the union over the immediate subpatterns\Added{, and on a
statement is the union over the branches of a conditional and empty on the statements not shown}.

\funccases{binds}
\end{auxfunction}

\begin{auxfunction}{\declaresS(t), \declaresS(\seq{\iota})}{sequence of variables}{aux:declares}
\added{Names declared anywhere in $t$, or bound by the imports of prefix $\seq{\iota}$, in textual order and with repeats, not descending into nested functions. It is empty on the statements not shown and on the empty prefix, and is the concatenation over the branches of a conditional and the cases of a match.}

\funccases{declares}
\end{auxfunction}

\begin{auxfunction}{\typeEntry{\Gamma}(q, \seq{\chi}, i), \typeEntry{\Gamma}(q, \seq{\chi})}{context entry, or sequence of context entries}{aux:type-entry}
Context entry associated with $i$-th name declared by mutual type region $\seq{\chi}$ of module $q$, under
context $\Gamma$: the class for a class declaration, and for a \pyinline{type} statement an alias carrying the type region
and $\Gamma$, under which the body of the type statement is resolved at each use (\ruleName{ty-alias}). For an entire type region $\seq{\chi}$, the sequence of these entries.

\funccases{type-entry}
\vspace{-\baselineskip}
\funccases{type-entry-region}
\end{auxfunction}

\begin{auxfunction}{\typeContext{\Gamma}(q, \seq{\chi}, i)}{context}{aux:type-context}
Context under which the type expressions of $\chi_i$, a definition of mutual type region $\seq{\chi}$ of module $q$, are
resolved, given context $\Gamma$ in which the type region appears: $\Gamma$ extended with the bindings of every
definition of $\seq{\chi}$ if $\chi_i$ is a \pyinline{type} statement, or of the definitions before $\chi_i$ if it is a
class definition, and then with the type parameters of $\chi_i$.

\funccases{type-context}
\end{auxfunction}

\begin{auxfunction}{\classTable{\Gamma}(q, \seq{\chi}, i), \classTable{\Gamma}(q, \seq{\chi})}{class table}{aux:class-table}
Class table contributed by the $i$-th declaration of mutual type region $\seq{\chi}$ of module $q$ under context
$\Gamma$, which defines the declared class for a class declaration and is empty for a \pyinline{type} statement. For an
entire type region $\seq{\chi}$, the class table defining every class declared in $\seq{\chi}$.

\funccases{class-table}
\vspace{-\baselineskip}
\funccases{class-table-region}
\end{auxfunction}

\begin{auxfunction}{\scopeS(Y, s)}{context}{aux:scope}
\Replaced{Context for the body $s$ of a function with parameters $Y$, or of a module with imported names $Y$, in which each variable declared or assigned in $s$ is unbound and the names in $Y$ keep their bindings. Defined if and only if the names declared in $s$, the names in $Y$ and the pattern variables of $s$ are pairwise distinct.}{Context in which each variable declared or assigned in $s$, the body of a scope, is unbound. Defined if and only if the names declared in $s$, the names in $Y$ and the pattern variables of $s$ are pairwise distinct, where $Y$ is the set of parameters of a function or of imported names of a module.} The context serves as an environment as well, $\Unbound$ being a value as well as an entry.

\funccases{scope}
\end{auxfunction}

\begin{auxfunction}{\captures(e)}{set of variables}{aux:captures-e}
Variables of the enclosing scope captured by lambdas within expression $e$; the generator variables of a
comprehension are local to it. On the cases not shown it recurses structurally, taking the
union over the immediate subexpressions.

\funccases{captures-e}
\end{auxfunction}

\begin{auxfunction}{\captures(\seq{g})}{set of variables}{aux:captures}
\replaced{Variables of the enclosing scope captured by lambdas within the qualifier sequence $\seq{g}$ of a comprehension.}{Variables of the enclosing scope captured by lambdas and nested function definitions within
$t$, or by lambdas within the qualifier sequence $\seq{g}$ of a comprehension. On the cases not shown it takes the union of $\capturesE$ over the subexpressions of $t$ and of
$\captures$ over its sub-statements, and on a module body agrees with $\assignsS$.}

\funccases{captures}
\end{auxfunction}

\begin{auxfunction}{\fv(e), \fv(\seq{e}), \fv(\tau)}{set of variables}{aux:fv}
Free variables of expression $e$, defined in the usual way, and of a sequence of expressions $\seq{e}$, as the union over the expressions; type variables of type $\tau$.

\end{auxfunction}

\section{Types}
\label{sec:types-auxiliary}

The functions on types and shapes introduced in \secref{types}, each defined by its cases.

\begin{auxfunction}{\baseType(\ell), \baseType(\tau)}{type}{aux:base-type}
Base type of literal $\ell$, and of type $\tau$ that type with a literal type replaced by the base type of
its literal.

\funccases{base-type}
\end{auxfunction}

\begin{auxfunction}{\joinTy{\Sigma}{\sigma}{\tau}, \bigJoinTy{\Sigma}{\seq{\tau}}}{type}{aux:join}
Larger of two comparable types, and their union otherwise. A sequence joins its first type with the join of
the rest, with \pyinline{Never} as the unit.

\funccases{join}
\end{auxfunction}

\begin{auxfunction}{\meetTy{\Sigma}{\sigma}{\tau}}{type}{aux:meet}
Cases after the first two apply only when $\sigma$ and $\tau$ are incomparable.

\funccases{meet}
\end{auxfunction}

\begin{auxfunction}{\subterms(\tau), \subterms(\seq{\tau})}{set of types}{aux:subterms}
Subterms of $\tau$ other than unions, including $\tau$ itself unless it is a union, and for a sequence
the subterms of each of its types. The candidates of Lemma~\ref{lem:candidate-instances} are joins of
these.

\funccases{subterms}
\end{auxfunction}

\begin{auxfunction}{\typeArgs{\Sigma}(\seq{\alpha}, \sigma, \tau), \typeArgs{\Sigma}(\seq{\alpha}, \seq{\sigma}, \seq{\tau})}{context binding type variables at types}{aux:type-args}
Type arguments for $\seq{\alpha}$ at which $\seq{\sigma}$ is instantiated to fit $\seq{\tau}$. Each type
argument is the join of the lower bounds on the type parameter, so the instantiation is the least under the lower
bounds; the type arguments are undefined if a type parameter does not occur in any pair. Single types $\sigma$ and
$\tau$ stand for sequences of length one.

\funccases{type-args}
\vspace{-\baselineskip}
\funccases{join-map}
\end{auxfunction}

\begin{auxfunction}{\elemTy{\Sigma}(\tau)}{type}{aux:elem-type}
Type of the elements drawn by a generator from a list, string, dictionary\replaced{, tuple or range type}{ or tuple type} or a union of these,
mirroring $\iterOp$ (\secref{auxiliary-values}). The element type of a tuple is the join of the base types of its
components. Undefined at any other type.

\funccases{elem-type}
\end{auxfunction}

\begin{auxfunction}{\shapes{\Sigma}(\tau, H), \shapes{\Sigma}(\seq{\tau})}{ordered set of shapes or of shape sequences}{aux:shapes}
Shapes of type $\tau$ that remain once the heads in $H$ are excluded. A union has no head to exclude, so
the disjuncts of the union are taken separately, each with the heads of $H$ typed at that disjunct. A
dictionary is refined by its keys rather than by a head, so a dictionary type gives the dictionary form
with no key bound. Every other type gives the rest form $\shapeAny{\tau}{H}$, except that no shape
remains when $H$ exhausts $\tau$: a literal type whose literal is in $H$, $\tyBool$ when both
$\pyinline{True}$ and $\pyinline{False}$ are in $H$, $\tyNone$ when $\pyinline{None}$ is in $H$, an instantiated class with an ancestor in $H$, and $\tyNever$. A sequence of types yields the sequences of shapes of its
components.

\funccases{shapes}
\end{auxfunction}

\begin{auxfunction}{\shapeTy(k)}{type}{aux:shape-type}
Least type of shape $k$, with $\Sigma \vdash k : \shapeTy(k)$ (\figref{shapes}): the class of a
constructor, componentwise for a tuple, and the type recorded by the form otherwise.

\funccases{shape-type}
\end{auxfunction}

\begin{auxfunction}{\matchedAs{\Sigma}(\Gamma, \tau, p)}{type}{aux:matched-as}
Type to which $\tau$ is narrowed if pattern $p$ matches: an over-approximation of the matched
values, as a type. Variable and wildcard patterns leave $\tau$ unchanged; literal patterns narrow to their literal
types; tuple patterns narrow each component; list and dictionary patterns narrow to the list or
dictionary type; constructor patterns narrow to the instance of their class at $\tau$, or to $\tau$ itself
when $\tau$ is below the instance; unions are narrowed disjunct by disjunct. Any other pattern does not
match any value of $\tau$, and narrows to $\tyNever$. The function emulates the narrowing of a match scrutinee
observed in mypy within a case, and gives the types of patterns named by \pyinline{as}
(\ruleName{check-as}).

\funccases{matched-as}
\end{auxfunction}

\begin{auxfunction}{\remaining{\Sigma}(\Gamma, \tau, p)}{type}{aux:remaining}
Type remaining from $\tau$ \Replaced{if pattern $p$ fails to match}{once pattern $p$ has failed to match}\Deleted{, or once each of the patterns $\seq{p}$ has
failed in turn}: an over-approximation of the unmatched values, as a type rather than a set of shapes.
Nothing remains after a variable or wildcard pattern, after a literal pattern at its own literal type or
at $\tyNone$, after the empty dictionary pattern at a dictionary type, after a tuple pattern each of whose
sub-patterns leaves nothing, or after a constructor pattern at a type below the instance of its class
each of whose sub-patterns leaves nothing at the field types\Replaced{. The}{; the} other Boolean literal remains after a
literal pattern at $\tyBool$; a tuple pattern with exactly one sub-pattern that leaves something narrows
that component; a union is narrowed \Replaced{disjunct by disjunct}{member by member}; and any other pattern leaves $\tau$ unchanged. The
function emulates the narrowing of a match scrutinee observed in mypy, so that a match which is exhaustive by the
remaining type is \Replaced{also exhaustive under}{exhaustive for} mypy. The remaining type is coarser than the residual
(\secref{well-formedness-patterns}), since \Replaced{constructor patterns with a refutable sub-pattern, and tuple patterns with two refutable components, leave the type untouched}{a constructor pattern with a refutable sub-pattern, or a
tuple pattern with two refutable components, leaves the type untouched}.

\funccases{remaining}
\end{auxfunction}

\begin{auxfunction}{\resolveOp{\Sigma}(\odot, \sigma_1, \sigma_2), \resolveOp{\Sigma}(\ominus, \sigma)}{type}{aux:resolve-op}
Result type of the most specific overload of an operator at the operand types. The overloads of $\odot$ at $(\sigma_1, \sigma_2)$
are the resolved overloads $((\tau_1, \tau_2), \tau)$ with $\Sigma \vdash \odot(\sigma_1, \sigma_2) : ((\tau_1, \tau_2), \tau)$
(\figref{operator-overloads}), ordered by componentwise subtyping in $\Sigma$ on their bounds $(\tau_1, \tau_2)$; the
overloads of $\ominus$ at $\sigma$ are defined in the same way from conclusions $\Sigma \vdash \ominus(\sigma) : (\tau, \tau')$.
\Replaced{The most specific overload is the least element under this order, so the function is undefined when no least element exists.
A least element exists when some overload applies and neither operand type is below \pyinline{Never}
(Lemma~\ref{lem:most-specific-overload}).}{The most specific overload, the least element under this order, exists when some overload applies and neither operand type is below
\pyinline{Never} (Lemma~\ref{lem:most-specific-overload}); otherwise the function is undefined.}

\funccases{resolve-op}
\end{auxfunction}

\section{Classes and modules}
\label{sec:auxiliary-classes}

The functions on module names, field lists and classes introduced in \secref{modules-programs}, and
the resolution of names in an environment.

\begin{auxfunction}{\submods(q)}{context}{aux:submods}
Context that binds each submodule of $q$ to a stub: a context entry $\bind{x}{\modStub{q.x}}$ for each module
$q.x$ of the program or of the predefined modules.

\funccases{submods}
\end{auxfunction}

\begin{auxfunction}{\containing(q, q')}{sequence of module names}{aux:containing}
Containing modules of $q'$, excluding $q$ and the modules containing $q$, in prefix order. A from-import of $q'$ in
module $q$ \replaced{checks}{loads} these as well as $q'$ itself.

\funccases{containing}
\end{auxfunction}

\begin{auxfunction}{\ownFields(\phi), \ownFields(\chi)}{sequence of bindings}{aux:own-fields}
Fields declared by a field list, or by a class definition, with their types, in order.

\funccases{own-fields}
\end{auxfunction}

\begin{auxfunction}{\typeParams{\Sigma}(C)}{sequence of type parameters}{aux:type-params}
Type parameters of class $C$, as declared.

\funccases{type-params}
\end{auxfunction}

\begin{auxfunction}{\base{\Sigma}(C)}{instantiated class or $\bot$}{aux:base}
Base class of $C$ at the type arguments \replaced{given in its definition, resolved under the context of the definition,}{given in its declaration, resolved under the context of the declaration,} or
$\bot$ if $C$ extends no class.

\funccases{base}
\end{auxfunction}

\begin{auxfunction}{\fields{\Sigma}(C)}{field scheme}{aux:fields}
Fields of class $C$ with their declared types, as a field scheme over the type parameters of $C$: the fields of
its base class, if any, instantiated at the type arguments of the base, followed by the fields declared by $C$, distinct
from the inherited fields by \ruleName{class-extend} (\figref{typing-top-level}). The fields of an instantiated class
$C[\seq{\tau}]$ are the instantiation $\tyInst{\fields{\Sigma}(C)}{\seq{\tau}}$.

\funccases{fields}
\end{auxfunction}

\begin{auxfunction}{\ancestors{\Sigma}(C)}{set of classes}{aux:ancestors}
\Replaced{Class $C$ together with all its base classes, as a scheme over the type parameters of $C$, each base
class instantiated at the type arguments given by its subclass. A class $D$ is among the ancestors, $D \in
\ancestors(\Sigma, C)$, when $D[\seq{\sigma}]$ occurs in the body for some $\seq{\sigma}$, which is all the evaluation
and shape rules test}{Class $C$ together with all its base classes}.

\funccases{ancestors}
\end{auxfunction}

\begin{auxfunction}{\fieldMap{\Sigma}(C, \seq{u})}{map from variables}{aux:field-map}
The argument supplied for each field of $C$: positional arguments fill the fields in order and keyword
arguments match by name. It is undefined unless the arguments cover the fields exactly, so constructor
patterns and expressions check and evaluate field by field without regard to the split of the arguments
between positional and keyword. On a class value, the same map, with the fields read from the value.

\funccases{field-map}
\end{auxfunction}

\begin{auxfunction}{\resolve(\rho, q)}{value}{aux:resolve}
Value that qualified name $q$ denotes in $\rho$, mirroring the static judgement $\Gamma \vdash q : \theta$
(\figref{resolve-name}): a simple name is looked up in $\rho$ directly, and a qualified name in the
environment of the loaded module that its prefix denotes. Undefined if the prefix denotes a module stub or a
value other than a module.

\funccases{resolve}
\end{auxfunction}

\begin{auxfunction}{\nameOf(q)}{string}{aux:name}
String literal naming module $q$, so that $\nameOf(\pymath{\_\_main\_\_}) = \pymath{"\_\_main\_\_"}$.

\end{auxfunction}

\section{Values}
\label{sec:auxiliary-values}

The functions on values used by the evaluation rules.
The comparisons \pyinline{==}, \pyinline{!=}, \pyinline{in} and \pyinline{not in} are given in terms of
structural equality $\eqOp{\Sigma}$ and membership $\contains{\Sigma}$. When defined, the
operators agree with Python, yielding $\aborts{\kappa}$ when Python raises an
exception. When undefined\paperonly{ (\secref{no-coercions})}, evaluation of the expression has no derivation, and an
implementation may abort or produce the result that Python gives (\paperonly{\secref{python-compatibility}}\speconly{\secref{conformance}}). Because
equality is partial, the order in which elements are compared is observable: equality and membership short-circuit
at the first element that decides the result, so an undefined comparison later in a sequence does not make
the whole comparison undefined.


\begin{auxfunction}{\hat{\odot}(v, v')}{result}{aux:binop}
Meaning of a strict binary operator.

\begin{center}
\begin{tabular}{>{\raggedright\arraybackslash}p{0.22\textwidth}>{\raggedright\arraybackslash}p{0.62\textwidth}}
\arrayrulecolor{gray!60}\hline
\rowcolor{grammarheadcolor}$\odot$ & $\hat{\odot}(v, v')$ \\
\pyinline{==} & $\val{\pyinline{False}}$ if both operands are \pyinline{nan}, and $\eqOp{\Sigma}(v, v')$ otherwise \\
\pyinline{!=} & $\val{\pyinline{True}}$ if both operands are \pyinline{nan} or $\eqOp{\Sigma}(v, v') = \val{\pyinline{False}}$, $\val{\pyinline{False}}$ if $\eqOp{\Sigma}(v, v') = \val{\pyinline{True}}$, and otherwise $\eqOp{\Sigma}(v, v')$ \\
\pyinline{in} & $\contains{\Sigma}(v', v)$ \\
\pyinline{not in} & $\val{\pyinline{True}}$ if $\contains{\Sigma}(v', v) = \val{\pyinline{False}}$, $\val{\pyinline{False}}$ if $\contains{\Sigma}(v', v) = \val{\pyinline{True}}$, and otherwise $\contains{\Sigma}(v', v)$ \\
\pyinline{<}\quad\pyinline{<=}\quad\pyinline{>}\quad\pyinline{>=} & \todo{Ordering, as in Python} \\
\pyinline{+}\quad\pyinline{-}\quad\pyinline{*}\quad\pyinline{/}\quad\pyinline{//}\quad\pyinline{\%}\quad\pyinline{**} & \todo{Arithmetic, as in Python. Float precision is left to the implementation. $\aborts{\errZeroDiv}$ where the divisor is zero, and $\aborts{\errType}$ where an operand is not a number}
\end{tabular}

\end{center}
\end{auxfunction}

\begin{auxfunction}{\hat{\ominus}(v)}{result}{aux:unop}
Meaning of a unary operator.

\begin{center}
\begin{tabular}{>{\raggedright\arraybackslash}p{0.22\textwidth}>{\raggedright\arraybackslash}p{0.62\textwidth}}
\arrayrulecolor{gray!60}\hline
\rowcolor{grammarheadcolor}$\ominus$ & $\hat{\ominus}(v)$ \\
\pyinline{not} & $\val{\pyinline{True}}$ if $v = \pyinline{False}$, and $\val{\pyinline{False}}$ if $v = \pyinline{True}$ \\
\pyinline{+} & $\val{v}$ if $v$ is a number, and $\aborts{\errType}$ otherwise \\
\pyinline{-} & $\val{-v}$ if $v$ is a number, and $\aborts{\errType}$ otherwise
\end{tabular}

\end{center}
\end{auxfunction}

\begin{auxfunction}{\eqOp{\Sigma}(v, v')}{result}{aux:eq}
Structural equality: literals by value, lists and tuples by their elements, dictionaries with the same
keys by their values, and objects of the same class by their field values, each through $\eqElems$.
\pyinline{None} compares with any value; otherwise equality is undefined between values of different \Replaced{forms, such as a number and a string, or a list and a tuple}{kinds}.
Dictionaries are compared in the entry order of the left operand.

\funccases{eq}
\end{auxfunction}

\begin{auxfunction}{\eqElems{\Sigma}(\seq{v}, \seq{v}\,')}{result}{aux:eq-elems}
Equality of two sequences of values: \pyinline{False} if the lengths differ, and otherwise compared from the first element, stopping at the first pair that is not equal, with that pair's result.

\funccases{eq-elems}
\end{auxfunction}

\begin{auxfunction}{\contains{\Sigma}(v, v')}{result}{aux:contains}
Whether $v'$ is an element of $v$: of a list or tuple by structural equality, of a dictionary as a key,
and of a string as a substring.

\funccases{contains}
\end{auxfunction}

\begin{auxfunction}{\containsElems{\Sigma}(\seq{v}, v)}{result}{aux:contains-elems}
Whether a sequence of values has an element structurally equal to $v$, taking the elements in order
and stopping at the first element that is not unequal to $v$, with that element's result.

\funccases{contains-elems}
\end{auxfunction}

\begin{auxfunction}{\elems(v)}{sequence of values}{aux:elems}
Elements of a list or tuple, or the characters of a string, in order.

\funccases{elems}
\end{auxfunction}

\begin{auxfunction}{\iterOp(v)}{sequence of values}{aux:iter}
Sequence drawn from $v$ by a generator: the keys when $v$ is a dictionary, the numbers $m, \ldots, n - 1$ when $v$ is a range with start $m$ and stop $n$, and otherwise the elements of $v$.

\funccases{iter}
\end{auxfunction}

\begin{auxfunction}{\getitem(v, v')}{result}{aux:getitem}
Subscripting: the value under key $v'$ of a dictionary, or the element at index $v'$ of a sequence,
counting from the end when the index is negative. Aborts with $\errKey$,
$\errIndex$ or $\errType$ when Python raises the exception of that name.

\funccases{getitem}
\end{auxfunction}

\begin{auxfunction}{\result(o)}{result}{aux:result}
Outcome of a statement read as the result of a call: a returned value, \pyinline{None} for a
statement that completes without returning, and an abort unchanged.

\funccases{result}
\end{auxfunction}

\begin{auxfunction}{\dispatch(\rho, v, \seq{p}, \seq{s})}{bindings and statement}{aux:dispatch}
The case taken by a \pyinline{match} statement for value $v$: the bindings and statement of the first pattern that
matches, or no bindings and \pyinline{pass} if no pattern matches.

\funccases{dispatch}
\end{auxfunction}

\begin{auxfunction}{\update(\delta, \delta')}{dictionary entries}{aux:update}
Entries $\delta$ with $w$ bound to $v$, replacing any earlier binding of $w$.

\funccases{update}
\end{auxfunction}

\begin{auxfunction}{\setitem(\delta, w, v)}{dictionary entries}{aux:setitem}
Dictionary entries $\delta$ with $w$ bound to $v$, replacing any earlier binding of $w$.

\funccases{setitem}
\end{auxfunction}
